% =============================================================================
%  preamble_stats_addon.tex
%  Insert this block into your TP4 preamble, immediately BEFORE \begin{document}.
%  It adds theorem environments + statistics macros. It does NOT redefine
%  anything from the `physics` package (\abs, \norm, \tr, \Tr, \pdv, \dv, ...),
%  so it coexists with your existing preamble unchanged.
% =============================================================================

% ---- Theorem-Umgebungen (gemeinsamer Zähler, nummeriert pro section) --------
% amsthm is already loaded in your TP4 preamble.
\theoremstyle{plain}
\newtheorem{satz}{Satz}[section]
\newtheorem{theorem}[satz]{Theorem}
\newtheorem{lemma}[satz]{Lemma}
\newtheorem{korollar}[satz]{Korollar}
\newtheorem{proposition}[satz]{Proposition}

\theoremstyle{definition}
\newtheorem{definition}[satz]{Definition}
\newtheorem{beispiel}[satz]{Beispiel}
\newtheorem{algorithmus}[satz]{Algorithmus}

\theoremstyle{remark}
\newtheorem{bemerkung}[satz]{Bemerkung}
\newtheorem*{bemerkung*}{Bemerkung}

% ---- Zahlmengen -------------------------------------------------------------
% \providecommand so this addon can be merged into a preamble that already
% defines some of these (e.g. \R, \C) without a "command already defined" abort.
% (In the current standalone main.tex nothing pre-exists, so these behave exactly
%  like \newcommand — zero change here, safety only on later merge.)
\providecommand{\R}{\mathbb{R}}
\providecommand{\N}{\mathbb{N}}
\providecommand{\Z}{\mathbb{Z}}
\providecommand{\Q}{\mathbb{Q}}
\providecommand{\C}{\mathbb{C}}
\providecommand{\indicator}{\mathbbm{1}}   % requires \usepackage{bbm} (no \mathbb{1} fallback: amsfonts blackboard has no digits)

% ---- Wahrscheinlichkeit & Statistik ----------------------------------------
% \ifdefined-guarded so merging into an existing preamble never aborts with
% "command already defined" (mirrors the guards on \rank/\diag/\sgn/\supp below).
\ifdefined\E\else\DeclareMathOperator{\E}{\mathbb{E}}\fi          % Erwartungswert
\ifdefined\Var\else\DeclareMathOperator{\Var}{Var}\fi
\ifdefined\Cov\else\DeclareMathOperator{\Cov}{Cov}\fi
\ifdefined\Corr\else\DeclareMathOperator{\Corr}{Corr}\fi
\ifdefined\Bias\else\DeclareMathOperator{\Bias}{Bias}\fi
\ifdefined\MSE\else\DeclareMathOperator{\MSE}{MSE}\fi
\ifdefined\Prob\else\DeclareMathOperator{\Prob}{\mathbb{P}}\fi    % \Pr ist LaTeX-intern -> \Prob
\ifdefined\argmax\else\DeclareMathOperator*{\argmax}{arg\,max}\fi
\ifdefined\argmin\else\DeclareMathOperator*{\argmin}{arg\,min}\fi
% Guarded: `physics` already defines \rank; guard the rest so the addon works
% whether or not physics is loaded.
\ifdefined\rank\else\DeclareMathOperator{\rank}{rank}\fi
\ifdefined\diag\else\DeclareMathOperator{\diag}{diag}\fi
\ifdefined\sgn\else \DeclareMathOperator{\sgn}{sgn}\fi
\ifdefined\supp\else\DeclareMathOperator{\supp}{supp}\fi

% bedingte Wahrscheinlichkeit / "gegeben": \Prob(A \given B).
% Robust with plain parens. For a stretchy bar use \left(\,\cdot\,\middle|\,\cdot\,\right).
\providecommand{\given}{\mid}
% Unabhängigkeit  X \indep Y
\providecommand{\indep}{\perp\!\!\!\perp}
% i.i.d.  X_1,\dots,X_n \iid \Normal(0,1)
\providecommand{\iid}{\overset{\text{iid}}{\sim}}
% konvergiert in Verteilung / Wahrscheinlichkeit
\providecommand{\convd}{\overset{d}{\to}}
\providecommand{\convp}{\overset{p}{\to}}

% ---- Verteilungen -----------------------------------------------------------
\providecommand{\Normal}{\mathcal{N}}
\providecommand{\Bernoulli}{\mathrm{Bernoulli}}
\providecommand{\Binomial}{\mathrm{Bin}}
\providecommand{\Poisson}{\mathrm{Poisson}}
\providecommand{\Uniform}{\mathrm{Unif}}
\providecommand{\Expo}{\mathrm{Exp}}

% ---- Vektoren ---------------------------------------------------------------
% EIN Makro, EINE Stelle zum Umschalten. \vek{x}, \vek{\beta}, \vek{1} (Einsvektor).
%
% Die Darstellung ist bewusst hier zentralisiert: die Vorlesung schreibt Vektoren
% an der Tafel unterstrichen, in einem gesetzten Skript ist das unüblich (Standard
% wäre fett). Diese Frage ist damit ein Einzeiler und keine Migration über 14.500
% Zeilen — vorher hing die Notation an ~390 einzelnen \underline{} im Fließtext.
%
% Umschalten auf die Buchkonvention: die beiden Zeilen tauschen. \boldsymbol (nicht
% \mathbf!) ist die richtige Wahl, weil \mathbf griechische Buchstaben NICHT fett
% setzt — \vek{\beta} und \vek{\mu} kommen im Skript vor.
% \providecommand{\vek}[1]{\underline{#1}}   % Tafelkonvention (wie im Video)
\providecommand{\vek}[1]{\boldsymbol{#1}}     % Buchkonvention (aktuell)

% ---- Schätzer / Transponieren ----------------------------------------------
\providecommand{\est}[1]{\hat{#1}}            % \est{\theta} -> \hat\theta
% \T beginnt mit ^ : Exponenten davor immer gruppieren, sonst "Double superscript"
% (also (X^{-1})\T statt X^{-1}\T). Nur im Mathe-Modus verwendbar.
\providecommand{\T}{^{\!\top}}                % A\T  (Transponierte)
